\exercisesection{3}

¶ 1. Let A be a ring, p a prime ideal of A, B a finitely generated A-algebra and let \(r_{1} \subset \cdots \subset r_{m}\) be a non-empty increasing sequence of ideals of B lying above p.~Show that there exist an element \(a \in A - p\) and a finite family \((y_{k})_{1 \leqslant k \leqslant q}\) of elements of B with the following properties:

\begin{enumerate}
\def\labelenumi{(\arabic{enumi})}
\item
  If we write \(\mathbf{B}' = \mathbf{A}[y_1, \ldots, y_q]\) and if \(\mathbf{S}\) is the set of powers of \(a\) , the ring \(\mathbf{S}^{-1}\mathbf{B}\) is integral over \(\mathbf{S}^{-1}\mathbf{B}'\) .
\item
  For all \(i\) such that \(1 \leqslant i \leqslant m\) there exists an integer \(h(i) \geqslant 0\) such that \(\mathfrak{r}_i\) contains \(y_1, \ldots, y_{h(i)}\) and for every polynomial \(\mathbf{P} \in \mathbf{A}[\mathbf{X}_{h(i)+1}, \ldots, \mathbf{X}_q]\) the relation \(\mathrm{P}(y_{h(i)+1}, \ldots, y_q) \in \mathfrak{r}_i\) implies that all the coefficients of \(\mathbf{P}\) belong to \(\mathfrak{p}\) .
\end{enumerate}

The ideal \(\mathfrak{r}_i\cap \mathbf{B}'\) is then generated by \(\mathfrak{p}\) and the \(y_{k}\) such that \(k\leqslant h(i)\) . (Argue by induction on \(m\) , thus reducing it to the case \(m = 1\) . Then let \(\mathfrak{r} = \mathfrak{r}_1\alpha\) ; argue by induction on the number \(n\) of generators of the A-algebra B. Let \((x_{i})_{1\leqslant i\leqslant n}\) be such a system of generators and let \(\mathfrak{s}\) be the ideal of \(\mathrm{A}[\mathrm{X}_1,\dots ,\mathrm{X}_n]\) consisting of those P such that \(\mathrm{P}(x_1,\ldots ,x_n)\in \mathfrak{r}\) ; then \(\mathfrak{s}\supset \mathfrak{pA}[\mathrm{X}_1,\ldots ,\mathrm{X}_n]\) ; distinguish two cases according to whether \(\mathfrak{s}\) is equal or not to \(\mathfrak{pA}[\mathrm{X}_1,\ldots ,\mathrm{X}_n]\) . In the second case, there is a polynomial \(\mathrm{Q}\in \mathrm{A}[\mathrm{X}_1,\ldots ,\mathrm{X}_n]\) none of whose coefficients is in \(\mathfrak{p}\) and which belongs to \(\mathfrak{s}\) . Show that it is possible to determine integers \(m_i\) ( \(2\leqslant i\leqslant n\) ) such that, if we write \(x_1' = \mathrm{Q}(x_1,\ldots ,x_n),x_i' = x_i - x_1^{m_i}\) ( \(2\leqslant i\leqslant n\) ), there exists an element \(b\in \mathrm{A} - \mathfrak{p}\) for which \(bx_{1}\) is integral over \(\mathrm{C} = \mathrm{A}[x_2',\ldots ,x_n']\) ; apply the induction hypothesis to the A-algebra

\[
\mathrm{B} _ {1} = \mathrm{A} [ x _ {2} ^ {\prime}, \dots , x _ {n} ^ {\prime} ]
\]

and the ideal \(\mathfrak{r} \cap \mathbf{B}_1\) .)

\begin{enumerate}
\def\labelenumi{\arabic{enumi}.}
\setcounter{enumi}{1}
\item
  Let K be a field and B the quotient of the polynomial ring K{[}X, Y{]} in two indeterminates by the ideal a of K{[}X, Y{]} generated by \(X^{2}\) and XY; denote by x and y the images in B of the elements X, Y; let A be the subring K{[}x{]} of B. Show that the elements x and \(y^{n}\) ( \(n \geqslant 0\) ) form a basis of B over K and that B is not integral over A. On the other hand, there exists in B no element which is algebraically free over A. Deduce a counter-example to Corollary 1 to Theorem 1 of no. 1 when A is not an integral domain.
\item
  Let K be an infinite field, L an extension field of K, \((x_{i})_{1 \leq i \leq n}\) a finite family of elements of L and d the transcendence degree of \(\mathrm{K}(x_{1}, \ldots, x_{n})\) over K. Show that there exist d elements \(z_{j} = \sum_{i=1}^{n} a_{ji} x_{i}\) of L (where the \(a_{ji} \in K\) ) such that \(\mathrm{K}[x_{1}, \ldots, x_{n}]\) is integral over \(\mathrm{K}[z_{1}, \ldots, z_{d}]\) . If we suppose further that \(\mathrm{K}(x_{1}, \ldots, x_{n})\) is a separable extension of K, show that the \(a_{ji}\) may be chosen such that the \(z_{j}\) form a separating transcendence basis of \(\mathrm{K}(x_{1}, \ldots, x_{n})\) over K (argue by induction on n when n \textgreater{} d).
\item
  Let B be an integral domain and A a subring of B. Suppose that there exist \(n \geqslant 1\) elements \(t_{1}, \ldots, t_{n}\) of B which are algebraically independent over A and such that B is integral over \(A[t_{1}, \ldots, t_{n}]\) . Let n be a maximal ideal of B and \(p = n \cap A\) . Then there is a prime ideal q of B contained in n, distinct from n and containing p.~(Reduce it first to the case where A is integrally closed by considering the integral closure \(A'\) of A (contained in the field of fractions of B) and the ring \(B' = B[A']\) ; in the case where A is integrally closed, use Theorem 3 of §2, no. 4.)
\end{enumerate}

¶ 5. Let A be an integral domain and B ⊃ A a subring of the field of fractions K of A; suppose that A is integrally closed in B and that B contains an element t such that B is a finitely generated A{[}t{]}-module.

\begin{enumerate}
\def\labelenumi{(\alph{enumi})}
\tightlist
\item
  Let F be a polynomial \(\neq 0\) of A{[}X{]} such that \(F(t)\) belongs to the conductor \(\mathfrak{f}\) of B with respect to A{[}t{]}; show that, if \(c\) is the dominant coefficient of F, \(c\) belongs to the radical of \(\mathfrak{f}\) . (By considering the rings A{[} \(c^{-1}\) {]} and B{[} \(c^{-1}\) {]}, reduce it to showing that, if \(c = 1\) , necessarily B = A{[}t{]}. On the other hand we may assume that F(t) \(\neq 0\) , otherwise \(t \in \mathbf{A}\) . For all \(x \in \mathbf{B}\) , by hypothesis
\end{enumerate}

\[
x \mathrm{F} (t) = \mathrm{F} (t) \mathrm{Q} (t) + \mathrm{R} (t)
\]

where Q and R are polynomials in A{[}X{]} and \(\deg(\mathrm{R}) < \deg(\mathrm{F})\) ; attention may be confined to the case where \(\mathrm{R}(t) \neq 0\) ; if \(y = \mathrm{R}(t)(\mathrm{F}(t))^{-1} = x - \mathrm{Q}(t)\) , show that t is integral over \(A[y^{-1}]\) and deduce that y is integral over A.)

\begin{enumerate}
\def\labelenumi{(\alph{enumi})}
\setcounter{enumi}{1}
\item
  Suppose further that A is Noetherian; let q be a prime ideal of B belonging to \(\operatorname{Ass}(B/\mathfrak{f})\) . Show that, if \(\omega\) is the canonical homomorphism from B to the field of fractions of \(B/q\) , \(\omega(t)\) is transcendental over the field of fractions of \(\mathbf{A}/(\mathfrak{q} \cap \mathbf{A})\) . (If \(p = q \cap A[t]\) , there is an element \(y \in B\) such that q is the set of \(z \in B\) for which \(zy \in \mathfrak{f}\) ; show that p is the conductor of \(B' = A[t] + By\) with respect to A{[}t{]}. Then argue by reductio ad absurdum using (a).)
\item
  Suppose that A is Noetherian. Let n be a maximal ideal of B and \(p = A \cap n\) . Show that, if \(A_{p} \neq B_{n}\) , there exists a prime ideal of B contained in n, distinct from n and containing p.~(Let \(\mathfrak{f}\) be the conductor of B with respect to A{[}t{]}. Suppose first that \(n \supset \mathfrak{f}\) ; consider a minimal element q in the set of prime ideals of B contained in n and containing \(\mathfrak{f}\); consider the canonical image of A in B/q and use (b) and Exercise 4. If \(n \not\supset \mathfrak{f}\), the local ring B\_n is equal to the local ring A{[}t{]}\_\{n\_1\}, where n\_1 = n ∩ A{[}t{]} (§ 1, Exercise 16). If t ∈ A\_p, show that A\_p = B\_n.~If on the contrary A\_p ≠ B\_n, deduce from § 2, Exercise 14 (c) that pA\_p{[}t{]} is a non-maximal prime ideal of A\_p{[}t{]} and conclude that pB\_n is a prime ideal of B\_n distinct from nB\_n.)
\end{enumerate}

\begin{enumerate}
\def\labelenumi{\arabic{enumi}.}
\setcounter{enumi}{5}
\tightlist
\item
  An integral domain A is called integrally Noetherian if it is Noetherian and if, for every finite sequence \((t_{i})_{1\leq i\leq n}\) of elements of an extension field of the field of fractions of A, the integral closure of \(A[t_{1},\ldots,t_{n}]\) is a finitely generated \(A[t_{1},\ldots,t_{n}]\) -module. Every finitely generated integral algebra over a field is an integrally Noetherian domain (no. 2, Theorem 2).
\end{enumerate}

\begin{enumerate}
\def\labelenumi{(\alph{enumi})}
\item
  Every finitely generated integral algebra over an integrally Noetherian domain is an integrally Noetherian domain. Every ring of fractions of an integrally Noetherian domain not reduced to 0 is an integrally Noetherian domain.
\item
  Let A be an integrally Noetherian domain \((t_{i})_{1\leq i\leq n}\) a finite sequence of elements of an extension field of the field of fractions K of A and B a subring of \(\mathbf{K}(t_{1},\ldots,t_{n})\) containing A and integral over A. Then B is a finitely generated A-module. (Note that the field of fractions of B is a finite algebraic extension of K, using Algebra, Chapter V, § 5, no. 7.)
\end{enumerate}

\begin{enumerate}
\def\labelenumi{\arabic{enumi}.}
\setcounter{enumi}{6}
\item
  Let A be an integrally Noetherian domain (Exercise 6), \((t_i)_{1 \leq i \leq n}\) a finite sequence of elements of an extension field of the field of fractions of A, B the ring \(A[t_1, \ldots, t_n]\) , \(n\) a maximal ideal of B, \(p = A \cap n\) , \(A'\) the integral closure of A in B and \(p' = n \cap A'\) . Then, if the local rings \(A_p'\) and \(B_n\) are distinct, there exists a prime ideal of B contained in \(n\) , distinct from \(n\) and containing \(p\) (``Zariski's Principal Theorem''). (Using Exercise 6 (b), reduce it to the case where \(A' = A\) , then reduce it to the case where A is a local ring with maximal ideal \(p\) . Show that \(B/n\) is then a finite algebraic extension of \(A/p\) (cf.~no. 4, Theorem 3); deduce that, for every ring C such that \(A \subset C \subset B\) , \(n \cap C\) is a maximal ideal of C. Then argue by induction on \(n\) , denoting by \(B_t\) the integral closure of \(A[t_1, \ldots, t_i]\) in B; distinguish two cases according to whether \(t_{i+1}\) is transcendental or algebraic over the field of fractions of \(B_t\) ; in the former case, use Exercise 4 and, in the latter, Exercise 5 (c).)
\item
  Let A be a ring. For A to be a Jacobson ring, it is necessary and sufficient that, for every maximal ideal \(m'\) of the polynomial ring A{[}X{]}, \(m' \cap A\) is a maximal ideal of A. (Note that if an integral domain B admits a Jacobson radical \(R \neq 0\) and if \(b \in R\) , \(b \neq 0\) , the intersection of B and a maximal ideal of B{[}X{]} containing \(1 + bX\) cannot contain b.)
\item
  Give an example of a Jacobson domain A and a Jacobson ring B containing A and such that there exists a maximal ideal of B whose intersection with A is not a maximal ideal.
\end{enumerate}

¶ 10. Let k be a field, L an infinite set and A the polynomial ring \(k[X_{\lambda}]_{\lambda \in L}\) . If \(k_{0}\) is the prime field contained in k, let b denote the cardinal of a transcendence basis of k over \(k_{0}\) and let us write \(c = \text{Card}(L)\) .

\begin{enumerate}
\def\labelenumi{(\alph{enumi})}
\item
  For every ideal a of A, show that there exists a system of generators S of a such that \(\operatorname{Card}(S) \leqslant c\) . (Consider the intersections of a and the subrings \(k[X_{\lambda}]_{\lambda \in J}\) for the finite subsets J of L.)
\item
  Suppose that \(c < b\) . Show that, if \(m\) is a maximal ideal of \(A\) , \(A / m\) is an algebraic extension of \(k\) . (Let \(K\) be the subfield of \(k\) generated over \(k_0\) by the coefficients of the polynomials forming a system of generators of \(m\) of cardinal \(\leqslant c\) ; let \(m_0 = m \cap K[X_\lambda]_{\lambda \in L}\) , so that \(m = m_0 \otimes_k k\) and
\end{enumerate}

\[
\mathrm{A} / \mathfrak {m} = (\mathrm{K} [ \mathrm{X} _ {\lambda} ] _ {\lambda \in \mathrm{L}} / \mathfrak {m} _ {0}) \otimes_ {\mathrm{K}} k.
\]

Note finally that there exists a K-isomorphism of the field of fractions of \(K[X_{\lambda}]_{\lambda \in L}/m_{0}\) onto an algebraic closure of k.) Deduce that A is then a Jacobson ring (use Exercise 8).

\begin{enumerate}
\def\labelenumi{(\alph{enumi})}
\setcounter{enumi}{2}
\tightlist
\item
  Suppose that \(b < c\) ; then \(\operatorname{Card}(k) \leqslant c\) , so that there exists a bijection \(\lambda \mapsto \xi_{\lambda}\) of a subset \(L'\) of \(L\) onto \(k\) , \(L'\) having a non-empty complement in \(L\) . Let \(\lambda_0 \in L - L'\) ; show that there exists a \(k\) -homomorphism \(u\) of \(A\) onto a subfield of the field of rational functions \(k(Y)\) such that \(u(X_{\lambda_0}) = Y\) ,
\end{enumerate}

\[
u (\mathbf {X} _ {\lambda}) = 1 / (\mathbf {Y} - \xi_ {\lambda})
\]

for \(\lambda \in \mathbf{L}'\) , \(u(\mathbf{X}_{\lambda}) = 0\) for \(\lambda \in \mathbf{L} - (\mathbf{L}' \cup \{\lambda_0\})\) ; if \(\mathfrak{m}\) is the kernel of \(u\) , \(\mathbf{A} / \mathfrak{m}\) is therefore not an algebraic extension of \(k\) .

\begin{enumerate}
\def\labelenumi{(\alph{enumi})}
\setcounter{enumi}{3}
\tightlist
\item
  Under the same hypotheses as in (c), show that A is not a Jacobson ring. (Note that there exist k-algebras B which are not Jacobson rings (for example local rings) such that \(\operatorname{Card}(B) = b\) .)
\end{enumerate}


